% AI-assisted manuscript source; responsible author: Carptopus.
\documentclass[11pt]{article}
\usepackage[a4paper,margin=28mm]{geometry}
\usepackage[T1]{fontenc}
\usepackage[utf8]{inputenc}
\usepackage{lmodern}
\usepackage{microtype}
\usepackage{amsmath,amssymb,amsthm,mathtools}
\usepackage{xcolor}
\usepackage[colorlinks=true,linkcolor=blue!55!black,citecolor=blue!55!black,urlcolor=blue!65!black]{hyperref}
\usepackage{fancyhdr}

\newtheorem{theorem}{Theorem}[section]
\newtheorem{lemma}[theorem]{Lemma}
\newtheorem{corollary}[theorem]{Corollary}
\theoremstyle{remark}
\newtheorem{remark}[theorem]{Remark}
\numberwithin{equation}{section}

\pagestyle{fancy}
\fancyhf{}
\fancyhead[L]{Carptopus}
\fancyhead[R]{Periodic infinite quaternion ringsets}
\fancyfoot[C]{\thepage}
\setlength{\headheight}{14pt}
\setlength{\parindent}{0pt}
\setlength{\parskip}{0.55em}
\setlength{\emergencystretch}{4em}

\title{Periodic infinite ringsets in the Lipschitz quaternions}
\author{Carptopus\\\href{mailto:carptopus@163.com}{carptopus@163.com}}
\date{23 August 2026}

\hypersetup{
  pdftitle={Periodic infinite ringsets in the Lipschitz quaternions},
  pdfauthor={Carptopus},
  pdfsubject={A local-global classification of periodic infinite ringsets in the Lipschitz quaternions},
  pdfkeywords={integer-valued polynomial, Lipschitz quaternion, ringset, null ideal, periodic word, local-global principle}
}

\begin{document}
\maketitle

\begin{center}
\small\itshape
Public Beta v0.1. The complete manuscript has passed independent proof and
manuscript-level zero-trust audits. External mathematical review and journal
peer review are pending.
\end{center}
\vspace{0.8em}
\begin{abstract}

Let $\mathbf L=\mathbb Z[\mathbf i,\mathbf j,\mathbf k]$ be the ring of
Lipschitz quaternions. Given a periodic word
$w:\mathbb Z\to\{\mathbf i,\mathbf j\}$, we study the infinite set

\[
S_w=\{a+w(a):a\in\mathbb Z\}\subset\mathbf L.
\]

A subset $S\subseteq\mathbf L$ is a ringset if the polynomials over the
rational quaternion algebra that are integer-valued on $S$ form a ring. We
give a complete criterion when a prescribed period $m$ is odd and squarefree.
The set $S_w$ is a ringset if and only if $w$ is nonconstant and, for every
prime $p\mid m$ with $p\equiv3\pmod4$, the restriction of $w$ to each residue
class modulo $p$ is nonconstant. Consequently, if every prime divisor of $m$
is congruent to $1$ modulo $4$, every nonconstant binary word of period $m$
produces an infinite ringset.

The proof combines the residue-null-ideal criterion for noncommutative
integer-valued polynomials with two local mechanisms. At split prime powers,
rank-one idempotents in
$\mathbf L/p^e\mathbf L\simeq M_2(\mathbb Z/p^e\mathbb Z)$ show that one
chosen letter above each real residue already forces a core set. At primes
congruent to $3$ modulo $4$, a Frobenius real-part extractor gives an explicit
right-multiplication obstruction whenever a residue fibre has only one
letter. Periodic congruence orbits and the Chinese remainder theorem then
give the global classification.

\end{abstract}
\section{Introduction}

Let $\mathbb D$ denote the rational Hamilton quaternion algebra and let

\[
\mathbf L
=\mathbb Z[\mathbf i,\mathbf j,\mathbf k],
\qquad
\mathbf i^2=\mathbf j^2=-1,
\qquad
\mathbf i\mathbf j=\mathbf k=-\mathbf j\mathbf i.
\]

The variable in $\mathbb D[x]$ is central, and polynomials are evaluated on
the right: if $f(x)=\sum_d c_dx^d$, then
$f(\alpha)=\sum_d c_d\alpha^d$. For $S\subseteq\mathbf L$, put

\[
\mathrm{Int}(S,\mathbf L)
=\{f\in\mathbb D[x]:f(S)\subseteq\mathbf L\}.
\]

This set is always a left $\mathbf L$-module, but in the noncommutative
setting it need not be closed under polynomial multiplication. Following
Werner \cite{werner2026}, $S$ is called a \emph{ringset} if
$\mathrm{Int}(S,\mathbf L)$ is a subring of $\mathbb D[x]$.

Werner \cite{werner2026} completely classified the finite ringsets in the Lipschitz and
Hurwitz quaternion rings. The infinite problem behaves differently. In
particular, a finite set can be decomposed into its minimal-polynomial
classes, whereas this reduction fails for infinite sets. Werner's
Proposition 5.6 gives a sufficient congruence-pairing condition and Example
5.9 constructs an infinite ringset whose individual minimal-polynomial
fibres are all non-ringsets.

We examine the simplest stationary version of that phenomenon. Let
$m\in\mathbb Z_{\ge1}$ and let
$w:\mathbb Z\to\{\mathbf i,\mathbf j\}$ satisfy $w(a+m)=w(a)$ for every
$a\in\mathbb Z$. The period $m$ is prescribed and need not be the least
period. Define

\[
S_w=\{a+w(a):a\in\mathbb Z\}.
\]

Every element $a+w(a)$ has minimal polynomial

\[
x^2-2ax+a^2+1.
\]

Thus distinct real parts lie in distinct minimal-polynomial classes, and
every such fibre is a noncentral singleton. Nevertheless, many of the
resulting infinite sets are ringsets.

Our main theorem is the following local--global classification.

\begin{theorem}

Let $m\ge1$ be odd and squarefree, and let
$w:\mathbb Z\to\{\mathbf i,\mathbf j\}$ have period $m$. Then $S_w$ is a
ringset of $\mathbf L$ if and only if both of the following hold:

\begin{enumerate}
\item $w$ is nonconstant;
\item for every prime $p\mid m$ with $p\equiv3\pmod4$, and every
$c\in\mathbb Z/p\mathbb Z$, the set

\[
\{w(a):a\equiv c\pmod p\}
\]

contains both $\mathbf i$ and $\mathbf j$.

\end{enumerate}
The criterion is independent of whether the displayed period is minimal.
At the split primes, no two-colour condition is needed.

\end{theorem}
\begin{corollary}

Suppose that $m$ is odd and squarefree and every prime divisor of $m$ is
congruent to $1$ modulo $4$. Then every nonconstant word of period $m$
produces a ringset $S_w$. In particular, among the $2^m$ labelled binary
words on $\mathbb Z/m\mathbb Z$, exactly $2^m-2$ satisfy the conclusion.
\end{corollary}

The finite-field core-set problem for $M_2(\mathbb F_q)$ has already been
classified algorithmically by Werner \cite{werner2022}. Rissner and Werner \cite{rissnerwerner} enumerated
core subsets within each similarity class, counted purely core subsets
globally, and obtained an asymptotic count for all core subsets. We do not
claim those local classifications and counts, the split quaternion
representation, or the residue-null-ideal bridge as new.
The contribution here is the explicit split prime-power lemma and its
combination with periodic congruence orbits and nonsplit obstructions to
obtain the above infinite-family if-and-only-if theorem.

\section{Null ideals and the residue criterion}

Let $A$ be a ring and $T\subseteq A$. Its null ideal is

\[
\mathcal N(T,A)
=\{f\in A[x]:f(t)=0\text{ for every }t\in T\}.
\]

It is always a left ideal. The set $T$ is called \emph{core} if
$\mathcal N(T,A)$ is a two-sided ideal. For $n\ge2$, write
$\overline S_n$ for the image of $S\subseteq\mathbf L$ in
$\mathbf L/n\mathbf L$.

We use the following consequence of Werner [4, Lemmas 5.3 and 5.4].

\begin{lemma}[residue criterion]

Let $S\subseteq\mathbf L$.

\begin{enumerate}
\item If $\overline S_n$ is core in $\mathbf L/n\mathbf L$ for every $n\ge2$,
then $S$ is a ringset.
\item Conversely, if for some $n$ there are
$g\in\mathcal N(\overline S_n,\mathbf L/n\mathbf L)$ and
$u\in\{\mathbf i,\mathbf j,\mathbf k\}$ such that
$gu\notin\mathcal N(\overline S_n,\mathbf L/n\mathbf L)$, then $S$ is
not a ringset.

\end{enumerate}
\end{lemma}
\begin{proof}

Every $f\in\mathbb D[x]$ can be written as $f=g/n$ with
$g\in\mathbf L[x]$ and $n\ge1$. By [4, Lemma 5.3], $f$ is integer-valued on
$S$ exactly when the reduction of $g$ belongs to
$\mathcal N(\overline S_n,\mathbf L/n\mathbf L)$.

If every residue image is core, reduction remains null after the coefficients
of $g$ are right-multiplied by $\mathbf i,\mathbf j$, or $\mathbf k$. Hence
$f\mathbf i,f\mathbf j,f\mathbf k$ are integer-valued on $S$. Since every
Lipschitz quaternion is an integral linear combination of
$1,\mathbf i,\mathbf j,\mathbf k$, [4, Lemma 5.4] gives multiplicative
closure. The same argument in reverse proves part 2 after lifting the
coefficients of the residue polynomial $g$ to $\mathbf L[x]$.
\end{proof}

The next elementary lemma records the paired-fibre calculation contained in
the proof of [4, Proposition 5.6]. We state it directly for a finite quotient;
we are not applying the global proposition itself to a finite ring.

\begin{lemma}[two letters over every real residue]

Let $q\ge2$, let $A=\mathbf L/q\mathbf L$, and let

\[
T=\{a+\mathbf i,a+\mathbf j:a\in\mathbb Z/q\mathbb Z\}.
\]

Then $T$ is core in $A$.

\end{lemma}
\begin{proof}

Take $f\in\mathcal N(T,A)$ and fix a real residue $a$. Divide $f$ by the
central monic polynomial

\[
\mu_a(x)=x^2-2ax+a^2+1.
\]

The remainder has the form $r(x)=\gamma_1x+\gamma_0$, and it vanishes at
$a+\mathbf i$ and $a+\mathbf j$. Subtracting gives

\[
\gamma_1(\mathbf i-\mathbf j)=0.
\]

Since

\[
(\mathbf i-\mathbf j)(-\mathbf i+\mathbf j)=2,
\]

we have $2\gamma_1=0$ in $A$; no invertibility of $2$ is used.

For $u\in\{\mathbf i,\mathbf j,\mathbf k\}$ and
$\alpha\in\{a+\mathbf i,a+\mathbf j\}$, direct right evaluation gives

\[
(ru)(\alpha)-r(\alpha)u
=\gamma_1(u\alpha-\alpha u).
\]

Every displayed commutator lies in $2A$, so the right-hand side is zero.
To pass from the remainder back to $f$, write $f=q_a\mu_a+r$. Since
$\mu_a$ has central coefficients,

\[
fu=(q_au)\mu_a+ru,
\qquad
((q_au)\mu_a)(\alpha)=0.
\]

Hence $(fu)(\alpha)=(ru)(\alpha)$ at both roots of $\mu_a$, and $fu$ is
null on $T$ for the three basis units. It follows by linearity
that the null ideal is stable under right multiplication by every constant in
$A$. If $h(x)=\sum_t r_tx^t\in A[x]$, then

\[
fh=\sum_t(fr_t)x^t=\sum_tx^t(fr_t)
\]

belongs to the null ideal because the latter is a left ideal and $x$ is
central. Hence the null ideal is two-sided.
\end{proof}

The same proof applies to any subset of the two-letter grid that contains both
letters above every real residue.

\section{A split prime-power lemma}

At a prime congruent to $1$ modulo $4$, the conclusion is much stronger: one
letter above each real residue is enough.

\begin{lemma}[split one-letter lemma]

Let $p\equiv1\pmod4$ be prime, let $e\ge1$, and put
$R=\mathbb Z/p^e\mathbb Z$ and $A=\mathbf L/p^e\mathbf L$. Suppose

\[
T\subseteq\{a+\mathbf i,a+\mathbf j:a\in R\}
\]

contains at least one point of real part $a$ for every $a\in R$. Then $T$ is
core in $A$.

\end{lemma}
\begin{proof}

Hensel lifting gives a unit $s\in R$ with $s^2=-1$. The assignments

\[
\mathbf i\longmapsto
\begin{pmatrix}s&0\\0&-s\end{pmatrix},
\qquad
\mathbf j\longmapsto
\begin{pmatrix}0&1\\-1&0\end{pmatrix}
\]

give an isomorphism $A\simeq M_2(R)$. Indeed,
$a+b\mathbf i+c\mathbf j+d\mathbf k$ maps to

\[
\begin{pmatrix}
a+bs&c+ds\\
-c+ds&a-bs
\end{pmatrix},
\]

and the four coordinates are recovered by dividing sums and differences by
the units $2$ and $s$.

For $u\in\{\mathbf i,\mathbf j\}$ define

\[
e_u^+=\frac{1+u/s}{2},
\qquad
e_u^-=\frac{1-u/s}{2}.
\]

These are complementary orthogonal idempotents. For every $a\in R$ and
$d\ge0$,

\[
(a+u)^d=(a+s)^d e_u^+ +(a-s)^d e_u^-.
\]

Under the matrix representation, their image modules are

\[
\begin{aligned}
\mathrm{im}(e_{\mathbf i}^+)&=R(1,0)^T,
&\mathrm{im}(e_{\mathbf i}^-)&=R(0,1)^T,\\
\mathrm{im}(e_{\mathbf j}^+)&=R(1,s)^T,
&\mathrm{im}(e_{\mathbf j}^-)&=R(1,-s)^T.
\end{aligned}
\]

For $u_1,u_2\in\{\mathbf i,\mathbf j\}$, place a generator of
$\mathrm{im}(e_{u_1}^+)$ beside a generator of
$\mathrm{im}(e_{u_2}^-)$. In the four cases

\[
(u_1,u_2)
=(\mathbf i,\mathbf i),(\mathbf i,\mathbf j),
(\mathbf j,\mathbf i),(\mathbf j,\mathbf j),
\]

the determinants may be chosen as

\[
1,\qquad -s,\qquad1,\qquad-2s,
\]

respectively. They are all units. Thus the two image modules are
complementary; if a matrix $X$ satisfies
$Xe_{u_1}^+=Xe_{u_2}^-=0$, then $X=0$.

Now take
$f(x)=\sum_dc_dx^d\in\mathcal N(T,A)$ and introduce the scalar-variable
matrix polynomial

\[
F(z)=\sum_dc_dz^d,
\qquad z\in R.
\]

For a fixed $z$, use the selected point of real part $z-s$ and multiply its
zero-evaluation equation on the right by its plus idempotent. This gives

\[
F(z)e_{u_1}^+=0.
\]

Similarly, the selected point of real part $z+s$ gives

\[
F(z)e_{u_2}^-=0.
\]

The complementarity just proved implies $F(z)=0$ for every $z\in R$. This
argument does not use uniqueness of polynomial functions over $R$.

Let $r\in A$ and let $fr$ denote uniform right multiplication of the
coefficients by $r$. Its scalar-variable polynomial is

\[
F_r(z)=\sum_d(c_dr)z^d=F(z)r=0.
\]

Applying the idempotent decomposition again shows that $fr$ vanishes at every
point of $T$. Hence the null ideal is stable under right multiplication by
all constants. As in Lemma 2.2, stability under arbitrary polynomial right
multiplication follows from the left-ideal property and centrality of $x$.
Therefore $T$ is core.
\end{proof}

The proof deliberately uses the matrix equations $Xe=0$ and the image modules
of the idempotents. No classification of left or right ideals is being
invoked.

\section{Periodic congruence fibres}

The arithmetic input is elementary but records exactly where squarefreeness
enters.

\begin{lemma}[periodic orbit lemma]

Let $w$ have period $m$, let $q\ge1$, and fix $a\in\mathbb Z/q\mathbb Z$.
The positions modulo $m$ attained by integers congruent to $a$ modulo $q$ are

\[
a+q\mathbb Z\pmod m
=\{r\in\mathbb Z/m\mathbb Z:r\equiv a
\pmod{\gcd(q,m)}\}.
\]

\end{lemma}
\begin{proof}

The subgroup generated by $q$ in $\mathbb Z/m\mathbb Z$ consists exactly of
the multiples of $\gcd(q,m)$. Translating by $a$ gives the formula.
\end{proof}

In particular, if $\gcd(q,m)=1$, every real residue fibre of $S_w$ contains
both letters exactly when $w$ is nonconstant. If $m$ is squarefree,
$p\mid m$, and $q=p^e$, then $\gcd(q,m)=p$, so a real residue fibre sees
exactly one position class modulo $p$.

The formula also shows that the criterion does not depend on choosing a
nonminimal period. If the least period is $d\mid m$ and a prime divisor of
$m$ does not divide $d$, each corresponding position class still projects
onto every position modulo $d$.

\section{Nonsplit missing-letter obstructions}

We next give the explicit obstruction used for necessity.

\begin{lemma}[missing-letter obstruction]

Let $p\equiv3\pmod4$ be prime and let

\[
T\subseteq\{a+\mathbf i,a+\mathbf j:a\in\mathbb F_p\}.
\]

If for some $a$ the fibre of real part $a$ is nonempty and contains only one
of the two letters, then $T$ is not core in $\mathbf L/p\mathbf L$.

\end{lemma}
\begin{proof}

In $(\mathbf L/p\mathbf L)[x]$ define the central polynomial

\[
\rho_p(x)=\frac{x+x^p}{2}.
\]

For $y=b+u$ with $u\in\{\mathbf i,\mathbf j\}$, the congruence
$p\equiv3\pmod4$ gives $u^p=-u$, and hence

\[
\rho_p(y)=b.
\]

The central Lagrange indicator of the real residue $a$ is

\[
\delta_a(x)=1-(\rho_p(x)-a)^{p-1}.
\]

It evaluates to $1$ above $a$ and to $0$ above every other real residue.

Suppose first that the fibre contains $a+\mathbf i$ but not
$a+\mathbf j$. Put

\[
h(x)=(x-(a+\mathbf i))\delta_a(x).
\]

The order of the two factors is intentional. Since $\delta_a$ has central
coefficients, $h$ vanishes on all of $T$. After uniformly right-multiplying
its coefficients by $\mathbf j$, evaluation at $\alpha=a+\mathbf i$ gives

\[
(h\mathbf j)(\alpha)
=\mathbf j\alpha-\alpha\mathbf j
=\mathbf j\mathbf i-\mathbf i\mathbf j
=-2\mathbf k\ne0.
\]

Thus the null ideal is not stable under right multiplication. If the fibre
contains only $a+\mathbf j$, use
$(x-(a+\mathbf j))\delta_a(x)$ and right multiplication by $\mathbf i$;
the value is $2\mathbf k\ne0$.
\end{proof}

Choosing integer representatives of the coefficients lifts the above null
polynomial to $\widehat h\in\mathbf L[x]$. Then
$\widehat h/p$ is integer-valued on every integer set whose residue image is
$T$, while the indicated right multiple is not. Lemma 2.1 therefore turns
the finite obstruction into a non-ringset certificate.

\section{Proof of the classification}

We now prove Theorem 1.1.

\paragraph{Sufficiency.}

Assume that $w$ is nonconstant and that every position class required in
condition 2 is bichromatic. We show that the image of $S_w$ is core modulo
every prime power $p^e$.

If $p\equiv1\pmod4$, every real residue modulo $p^e$ is represented by at
least one element of $S_w$. Lemma 3.1 applies without any two-colour
condition.

Suppose $p=2$. Since $m$ is odd,
$\gcd(2^e,m)=1$. Lemma 4.1 and nonconstancy show that every real residue
fibre modulo $2^e$ contains both letters. Lemma 2.2 applies.

Now let $p\equiv3\pmod4$. If $p\nmid m$, then again
$\gcd(p^e,m)=1$, and Lemmas 4.1 and 2.2 apply. If $p\mid m$, squarefreeness
gives $\gcd(p^e,m)=p$. Each real residue fibre therefore sees exactly one
position class modulo $p$, which is bichromatic by condition 2. Lemma 2.2
again applies.

It remains to pass from prime powers to arbitrary denominators. Write
$n=\prod_rp_r^{e_r}$. The Chinese remainder theorem gives

\[
\mathbf L/n\mathbf L
\simeq\prod_r\mathbf L/p_r^{e_r}\mathbf L.
\]

For $A=\prod_rA_r$ and an arbitrary subset $T\subseteq A$, polynomial
coefficients and right evaluation are componentwise. Therefore

\[
\mathcal N(T,A)
\simeq
\prod_r\mathcal N(\pi_rT,A_r).
\]

This identity does not require $T$ to be the Cartesian product of its
projections: a component polynomial vanishes on every element of $T$ exactly
when it vanishes on the corresponding projection. Since every prime-power
projection of $S_w$ is core, its image modulo $n$ is core. Lemma 2.1 now
shows that $S_w$ is a ringset.

\paragraph{Necessity.}

If $w$ is constant, reduce modulo $3$. Every real residue fibre then contains
only one letter, so Lemma 5.1 and Lemma 2.1 show that $S_w$ is not a ringset.

Next suppose that $p\mid m$, $p\equiv3\pmod4$, and some position class
$c\pmod p$ is monochromatic. Fix a real residue $a\equiv c\pmod p$. By
Lemma 4.1, the fibre of the image of $S_w$ above $a$ sees precisely that
position class, and hence contains only one letter. Lemma 5.1 again gives a
null polynomial whose right multiple is non-null. Lemma 2.1 implies that
$S_w$ is not a ringset. This proves both necessary conditions and completes
the theorem.

\section{Examples and the squarefree boundary}

If $m=5$, the only prime divisor is split. Hence every one of the
$2^5-2=30$ nonconstant labelled words produces a ringset. More generally,
Corollary 1.2 gives $2^m-2$ examples whenever all prime divisors of $m$ are
$1$ modulo $4$.

For $m=3$ or $m=7$, condition 2 is impossible: a position class modulo $m$
contains only one position. Thus no word of these prescribed periods gives a
ringset. Composite periods can behave differently. For $m=15$, the condition
requires each of the three position classes modulo $3$ to be bichromatic. For
$m=21$, both every class modulo $3$ and every class modulo $7$ must be
bichromatic. The length-$m$ word obtained by alternating the two letters on
the standard representatives $0,1,\ldots,m-1$ satisfies these requirements
for $m=15,21$, and similarly for $m=35$.

Squarefreeness is used only in the equality
$\gcd(p^e,m)=p$ when a nonsplit prime divides $m$. If $p^2\mid m$, a fibre
modulo $p^e$ may see a position class modulo $p^2$ or a higher power. Being
bichromatic modulo $p$ no longer guarantees a paired fine fibre, while a
missing letter in one fine fibre is not automatically detected by the
finite-field polynomial of Lemma 5.1. A classification for nonsquarefree
periods therefore requires genuinely local-ring information about core
subsets of $\mathbf L/p^e\mathbf L$. We make no claim about that boundary
here.

\section{Computational calibration}

The proof above is symbolic and does not depend on computation. During
discovery, the degree-$\le2p$ truncated evaluation kernel over $\mathbb F_p$
was checked for all three choices \texttt{i only}, \texttt{j only}, or \texttt{both} over every
real residue:

\begin{itemize}
\item all $3^3=27$ patterns for $p=3$;
\item all $3^5=243$ patterns for $p=5$;
\item all $3^7=2187$ patterns for $p=7$;
\item representative patterns for $p=11$ and $p=13$.

\end{itemize}
Within this degree-$\le2p$ truncation, for $p\equiv3\pmod4$ only the
all-bichromatic pattern was stable; for $p\equiv1\pmod4$, every pattern was
stable. These checks motivated the local lemmas but do not prove them or
establish stability of the full null ideal. The discovery program is retained
with the research record and explicitly labels itself as finite-field
calibration.

AI assistance was used to help formulate finite test cases, review the
exploratory program, and compare its reported outputs with the symbolic
argument. The author executed and inspected the computations; no computed
case is used as a proof of the theorem.

\section{Prior-work boundary}

The following components are prior work and are not claimed here:

\begin{itemize}
\item the general connection between ringsets and null-ideal sets \cite{frisch2018};
\item the algorithmic classification of finite core subsets of
$M_2(\mathbb F_q)$ \cite{werner2022}, and the within-class, purely-core, and asymptotic
counts in \cite{rissnerwerner};
\item the finite Lipschitz/Hurwitz ringset classification, the residue criterion,
the right-multiplication criterion, and the infinite sufficient condition
\cite{werner2026}.

\end{itemize}
As of the literature search completed on 23 August 2026, we found no public
source directly stating the odd-squarefree periodic criterion of Theorem 1.1.
This search result is a bounded prior-work assessment, not a claim of absolute
global priority. The result has not yet received external expert review or
formal peer review.

\section*{Declaration of generative AI and AI-assisted technologies}

During the research and preparation of this work, the author used OpenAI
Codex to assist with literature searches, symbolic exploration, exploratory
code review, proof checking, and language drafting. The author reviewed and
edited all generated output and takes full responsibility for the content of
this manuscript.


\begin{thebibliography}{9}

\bibitem{frisch2018}
Sophie Frisch,
\emph{Polynomial functions on subsets of non-commutative rings---a link between ringsets and null-ideal sets},
ITM Web of Conferences \textbf{20} (2018), 01003.
\href{https://doi.org/10.1051/itmconf/20182001003}{doi:10.1051/itmconf/20182001003}.

\bibitem{rissnerwerner}
Roswitha Rissner and Nicholas J. Werner,
\emph{Counting core sets in matrix rings over finite fields},
Linear Algebra and its Applications \textbf{720} (2025), 1--25.
\href{https://doi.org/10.1016/j.laa.2025.04.006}{doi:10.1016/j.laa.2025.04.006};
\href{https://arxiv.org/abs/2405.04106}{arXiv:2405.04106v2}.

\bibitem{werner2022}
Nicholas J. Werner,
\emph{Null ideals of subsets of matrix rings over fields},
Linear Algebra and its Applications \textbf{642} (2022), 50--72.
\href{https://doi.org/10.1016/j.laa.2022.02.007}{doi:10.1016/j.laa.2022.02.007}.

\bibitem{werner2026}
Nicholas J. Werner,
\emph{Integer-valued polynomials on subsets of quaternion algebras},
Journal of Algebra \textbf{686} (2026), 195--219.
\href{https://doi.org/10.1016/j.jalgebra.2025.08.011}{doi:10.1016/j.jalgebra.2025.08.011};
\href{https://arxiv.org/abs/2412.20609}{arXiv:2412.20609v2}.

\end{thebibliography}

\end{document}
